% !TEX root = ../../UPDATED_main.tex
\subsection*{Original SatuTe defines the branch-specific test}

SatuTe tests whether the subtrees separated by a focal branch retain detectable statistical dependence. For the slowest non-stationary component, original SatuTe averages the coherence coefficient in \cref{eq-coherence} across \(n\) independently evolving sites and standardizes that mean with its estimated null variance. The resulting one-sided statistic rejects subtree independence when it exceeds the selected normal threshold. When the tree is inferred from the same alignment, SatuTe applies a taxon-pair Bonferroni threshold to account for topology selection.

This branch, null hypothesis, and decision rule define the primary model used here. The extensions below alter how coherence is summarized under models with distinct non-stationary eigenvalues or explicit invariant sites; they do not change the biological question posed by SatuTe.

\subsection*{A weighted coefficient extends the SatuTe statistic}

The spectral decomposition underlying SatuTe gives
\begin{equation}
\frac{P(\partial\mid t)}{P(\partial_A)P(\partial_B)}
=
1+\sum_{i=1}^{d} C_{\partial i}\exp(\lambda_i t).
\label{eq-central}
\end{equation}
Here \(d\) is the number of non-stationary eigenvalues: \(d=3\) for DNA and \(d=19\) for an amino-acid model. The coefficient \(C_{\partial i}\) measures agreement between the two subtrees in spectral direction \(h_i\), whereas \(\exp(\lambda_i t)\) measures how much of that direction persists across the focal branch. Because all non-stationary eigenvalues are negative, components with more negative \(\lambda_i\) decay faster.

For pattern \(\partial\), we define the eigenvalue-weighted coherence as
\begin{equation}
C_{\partial}^{\lambda}(t)
=
\sum_{i=1}^{d} w_i(t)C_{\partial i},
\qquad
w_i(t)=\exp\{(\lambda_i-\lambda_1)t\}.
\label{eq-weighted-coherence}
\end{equation}
The relative weight \(w_i(t)\) sets \(w_1(t)=1\) and downweights components that decay faster than the slowest non-stationary component. Original SatuTe therefore remains the reference scale, while additional components contribute according to the fitted branch length.

For an alignment of \(n\) independently evolving sites, the empirical coefficient is
\begin{equation}
\hat C_{\lambda}(t)=\frac{1}{n}\sum_{s=1}^{n}C_{\partial_s}^{\lambda}(t).
\label{eq-weighted-mean}
\end{equation}
Including the mode weights in the original variance estimator gives
\begin{equation}
\hat\sigma_{\lambda}^{2}(t)
=
\sum_{i=1}^{d}\sum_{j=1}^{d}
w_i(t)w_j(t)\hat\sigma_{A,i,j}^{2}\hat\sigma_{B,i,j}^{2}.
\label{eq-weighted-variance}
\end{equation}
The corresponding statistic is
\begin{equation}
Z_{\lambda}(t)=\frac{\sqrt{n}\hat C_{\lambda}(t)}{\hat\sigma_{\lambda}(t)}.
\label{eq-weighted-z}
\end{equation}
The original decision rule is unchanged: the focal branch is called phylogenetically informative when \(Z_{\lambda}(t)>z_\alpha\), with threshold \(z_{\alpha/(m_A m_B)}\) when the tested topology is inferred from the same alignment.

\subsection*{The Jukes--Cantor formulation is unchanged}

The JC model has one stationary eigenvalue and a three-dimensional non-stationary eigenspace with \(\lambda_1=\lambda_2=\lambda_3\). Therefore \(w_i(t)=1\) for all non-stationary modes in \cref{eq-weighted-coherence}, and
\begin{equation}
C_{\partial}^{\lambda}(t)
=
\sum_{k=1}^{3} C_{\partial 1,k}
=
C_{\partial 1},
\label{eq-jc-collapse}
\end{equation}
which is the multiplicity-three expression used by original SatuTe. The weighted statistic thus makes no change when the substitution model does not distinguish among non-stationary directions.

\subsection*{Invariant categories require a corrected null likelihood}

An explicit invariant category is a boundary case of a rate mixture. Let \(q_c\) and \(r_c\) be the proportion and rate of category \(c\), and let \(A_{sc}\) and \(B_{sc}\) be the two subtree-pattern probabilities for site pattern \(s\). For a positive-rate category, the infinite-branch null contribution is \(q_cA_{sc}B_{sc}\). For the invariant category \(c=0\), however, \(r_0=0\), and the transition matrix remains the identity even as the focal branch length tends to infinity. Its null contribution is therefore the joint probability
\begin{equation}
J_{s0}
=
\sum_a \pi_a L_{A,s0}(a)L_{B,s0}(a),
\label{eq-invariant-joint}
\end{equation}
rather than the product \(A_{s0}B_{s0}\). The corrected mixture null is
\begin{equation}
p_\infty(s)
=
q_0J_{s0}
+
\sum_{c:r_c>0}q_cA_{sc}B_{sc}.
\label{eq-invariant-null}
\end{equation}

For the positive-rate categories, define the model-weighted site score
\begin{equation}
u_s(t)
=
\frac{
\displaystyle
\sum_{c:r_c>0}q_cA_{sc}B_{sc}
\sum_{i=1}^{d}e^{\lambda_i r_ct}C_{sc,i}
}{p_\infty(s)}.
\label{eq-soft-mixture-score}
\end{equation}
The invariant component enters only the denominator and contributes no finite-versus-infinite-branch signal, giving the exact identity
\begin{equation}
\frac{p_t(s)}{p_\infty(s)}=1+u_s(t).
\label{eq-soft-mixture-identity}
\end{equation}
In computation, all positive-rate persistence weights are divided by the common factor \(e^\delta\), where \(\delta=\max_{c:r_c>0,i}\lambda_i r_ct\). This prevents underflow without changing the standardized statistic. The refined invariant-null test uses the sample mean and variance of these globally rescaled site scores.

A counterfactual construction that replaces \(J_{s0}\) with \(A_{s0}B_{s0}\) treats the invariant category as if it became independent across an infinite focal branch. It does not satisfy \cref{eq-soft-mixture-identity} under a \(+I\) model and need not have null expectation zero.
